\begin{textsample}{POS Dim 4 – ac – Score 10.00 – ac/ac\_ef\_2.txt}  \label{ex:f4_pos_261}
3.
PARTE DIVERSIFICADA E ESPECIFICIDADES DO ESTADO DO ACRE O currículo de Arte do Estado do Acre propõe que cada uma das dimensões de conhecimento seja desenvolvida, sem perder de vista a especificidade regional do Estado.
Ele está inserido dentro da maior \textbf{floresta} tropical do mundo, cuja \textbf{população} constitui -se pelas mais diversas origens, desde os povos originários, ameríndios, passando pelas ondas de \textbf{migrações} dos ciclos da borracha.
Estes ciclos trouxeram grande influência da cultura nordestina para a região, chegando aos fluxos \textbf{migratórios} mais recentes, trazendo uma grande \textbf{população} do \textbf{sul} e sudeste do Brasil.
Ao buscar refletir e valorizar tal diversidade, este currículo oferece sugestões de trabalhos com aspectos da cultura tradicional e popular, valorizando as festas populares como a Cavalhada ( Sena Madureira ), Marujada ( Rio Branco e Cruzeiro do Sul ), Novenário Nossa Senhora da Glória ( Cruzeiro do Sul ), Reisado ( Tarauacá ), Pastorinhas ( Rio Branco e Xapuri ), Festa Junina ( Presente em todo estado ), Festas do Daime, Festivais \textbf{Indígenas} ( MaririYanawa, Festa Huni-Kuin ), Festival do Açaí ( Feijó ), Festival do Peixe ( Sena Madureira ), Festa de São Sebastião ( Xapuri ), Expoacre ( Rio Branco ), ExpoJuruá ( Cruzeiro do Sul ), entre outras.
Nestas festividades fica clara a diversidade presente no Estado.
Não se restringindo a tais festas, o professor deve se relacionar, em cada município, com as manifestações locais e regionais, valorizando -as como produção de cultura.
Destaca -se a importância de a aula de Arte contemplar temas das culturas \textbf{indígena} e afro-brasileira, cumprindo as disposições legais e valorizando as origens da \textbf{população} local.
O professor de Arte deve, ainda, contemplar em seu planejamento, a pesquisa sobre \textbf{artistas} locais e equipamentos culturais de cada município, promovendo visitas, entrevistas e pesquisas in loco, considerando que a arte se dá na prática social da comunidade. 4.
ORIENTAÇÕES DE APLICABILIDADE DO COMPONENTE ARTE A Arte na escola tem por desiderato proporcionar ao aluno a ampliação do acesso a experiências estéticas, onde o aluno torna -se protagonista desse processo de aprendizagem.
Na oportunidade, o aprendente expressa seus sentimentos e desenvolve sua criatividade.
Em suma, em consonância com a BNCC, o componente curricular apresenta propostas que tem como cerne as quatro linguagens artísticas ( Artes Visuais, Dança, Música e Teatro ), agora denominadas de Unidades Temáticas, acrescidas de outra, chamada Artes Integradas ( uma forma de integração das linguagens e suas diferentes práticas, possibilidades e tecnologias ).
Ademais, soma -se a esse pressuposto, a importância de associação das Unidades Temáticas às seis Dimensões de Conhecimento ( Criação ; Crítica ; Fruição ; Estesia ; Expressão e Reflexão ), no intuito de garantir os direitos de aprendizagens dos alunos.
Tomando este propósito por referência, seguem os objetivos sugeridos para as linguagens artísticas dos Ensinos Fundamental I e II :

% matched lemmas: artista, floresta, indígena, migratório, migração, população, sul
\end{textsample}
